\section{Related Work}
\label{sec:related-work}

% Style note (v2, 2026-05-07): compressed from ~700 to ~250 words to fit
% the 4-page short-paper budget per plan v4. Five paragraphs reduced to
% three; less-core comparators (RankSteer, REALM, RefRank, FIRST,
% self-calibrated listwise) folded into single sentences or dropped.

\paragraph{LLM reranking paradigms.}
LLM zero-shot reranking is organized around pointwise, pairwise, listwise,
and setwise prompting.
Pointwise rerankers score each query--document pair
independently~\citep{nogueira2020monot5,sachan2022upr,ma2024rankllama};
pairwise methods choose between two documents and aggregate local
preferences~\citep{qin2024prp,luo2024prpgraph}; listwise methods prompt the
LLM to emit a permutation over a candidate
block~\citep{sun2023rankgpt,ma2023zeroshotlistwise,pradeep2023rankvicuna,
pradeep2023rankzephyr,zhang2025rankwithoutgpt}.
Setwise prompting~\citep{zhuang2024setwise} occupies a middle ground: each
call observes more than two documents, but the output is a short label
rather than a full permutation. MaxContext belongs to this family. It
keeps the short-label setwise interface but changes the granularity of
each comparison from a small local window to the whole current live pool.

\paragraph{Setwise lineage and aggregation alternatives.}
\citet{zhuang2024setwise} introduce setwise prompting with a small window
plus heapsort or bubblesort; we build on the public
\texttt{llm-rankers} implementation~\citep{ielab_llmrankers} but replace
its fixed letter labels with numeric labels so that the whole live pool is
addressable. Recent setwise extensions warm-start from
BM25~\citep{podolak2025setwiseinsertion} or train a setwise reasoner with
reinforcement learning~\citep{zhuang2025rankr1}. A separate line of work
extracts more global information from $k$-wise judgments by altering the
schedule or aggregation rule: TourRank uses a points-based
tournament~\citep{chen2025tourrank}, BLITZRANK builds preference graphs
from $k$-wise comparisons~\citep{agrawal2026blitzrank}, and PRP-Graph
aggregates pairwise units with target-label
probabilities~\citep{luo2024prpgraph}. MaxContext is orthogonal: it does
not warm-start, train, or aggregate; it asks what setwise selection looks
like when the candidate window is the live pool itself.

\paragraph{Long-context order, position bias, and efficiency accounting.}
Whole-pool prompting does not imply order invariance. Long-context models
under-use information at some prompt
positions~\citep{liu2024lostmiddle,tang2024foundmiddle}, and label-based
LLM judgments are sensitive to option labels and presentation
order~\citep{zheng2024mcqbias,shi2025positionbias,pradeep2023rankzephyr}.
We therefore include explicit input-order controls rather than assume
whole-pool visibility makes MaxContext robust. We also separate efficiency
axes: prior work emphasizes that LLM reranker cost is not summarizable
by a single proxy because calls, tokens, FLOPs, and wall-clock can move
in different
directions~\citep{zhuang2024setwise,podolak2025setwiseinsertion,
peng2025flopsreranking}; we report each, and keep the central MaxContext
claim narrowly framed in serial LLM calls and observed wall-clock.
Statistical comparisons follow standard IR
practice~\citep{smucker2007comparison,urbano2019significance}.
